\subsection{Ordinary arithmetic over every explicit global function field}
\label{C:global-function-arithmetic-subsection}

Here $q$ is an odd prime power and a finite extension of
$\mathbf F_q(t)$ is supplied as a finite tower of monic
irreducible polynomial presentations. Inseparable steps are
allowed. An algorithm means a terminating ordinary algorithm;
no running-time bound is asserted. Finite searches below are
dovetailed when their individual termination is not known.
The use of global arithmetic in this subsection is limited to
the Albert--Brauer--Hasse--Noether theorem and equality of
period and index over global fields, in their full
characteristic-positive versions.

\begin{lemma}[Proved: field arithmetic and enumeration of all completions]
\label{C:global-function-field-arithmetic}
Uniformly for every supplied finite tower $K/\mathbf F_q(t)$,
there are ordinary algorithms for field arithmetic, equality,
enumeration, linear algebra, and univariate factorization.
The same assertions hold after every further supplied finite
tower, including inseparable towers.

There is an ordinary enumeration, with possible repetitions,
of all places $v$ of $K$. Each enumerated place includes a
finite field $\kappa_v$, a named uniformizer $u_v\in K$, and
an explicit embedding
\[
             K\longrightarrow \kappa_v((U)),
             \qquad u_v\longmapsto U,
\]
whose completion is an isomorphism
$K_v\cong\kappa_v((U))$.
For a named nonzero element of $K$, its valuation, leading
residue, and any prescribed finite Laurent expansion are
computable. All rings
$\mathcal O_v/(u_v^N)\cong\kappa_v[U]/(U^N)$ are consequently
effective finite rings. An element of $K$ with any prescribed
finite Laurent approximation in this completion can be found
by an ordinary algorithm.
\end{lemma}
\begin{proof}
The power bases of the supplied tower give exact operations
and linear algebra. Rational-function names and finite tower
coordinates give an effective countable enumeration; equality
removes repetitions if desired.

\emph{Factorization, including inseparability.}
For a polynomial over $\mathbf F_q(t)$, clear denominators
and take a primitive polynomial $F(t,X)\in\mathbf F_q[t,X]$.
Every factor has total degree at most the total degree of $F$.
There are only finitely many polynomials over $\mathbf F_q$
within that bound. Exact division and repeated proper
factorization therefore give its irreducible factors.
Factors depending only on $t$ become units; Gauss's lemma
gives the desired factorization over $\mathbf F_q(t)$.

The Frobenius-root proof of
Lemma~\ref{C:function-field-pth-roots} applies with
$C=\mathbf F_q$: inverse Frobenius on this finite field is
effective. Explicitly, put $k=\mathbf F_q(t)$ and
$k^p=\mathbf F_q(t^p)$. For a known $k$-basis
$e_1,\ldots,e_s$ of $K$, the elements
$t^a e_i$, $0\leq a<p$, form a $k^p$-basis of $K$, whereas
the $e_i^p$ form a $k^p$-basis of $K^p$.
Solving for an input in their span decides membership in
$K^p$ and computes the unique root. This does not require
$K/k$ to be separable.

The entire tower-induction proof of
Theorem~\ref{C:function-field-factorization} now applies.
For clarity, its hypotheses and termination arguments are
as follows. Over a separable simple extension $E(\alpha)/E$,
shift a separable input $f(X)$ to $f(X-c\alpha)$ and enumerate
$c\in E$ until its norm to $E[X]$ is separable.
Only finitely many $c$ cause collisions among its conjugate
roots, and $E$ is infinite. Factor the norm over $E$ and
take gcds over $E(\alpha)$ to recover the factors. A general
simple extension first has a separable intermediate field
$E(\alpha^{p^a})$, followed by a purely inseparable extension.
Raise the polynomial to its $p^a$-th power to descend its
coefficients, factor there, and take all possible polynomial
$p$-th roots using the preceding root test. Each descended
irreducible factor has one irreducible factor after purely
inseparable extension, with multiplicity a power of $p$;
the proof in the cited theorem establishes this without
assuming perfectness. Repeated gcds with the derivative and
the reduction $f(X)=h(X^p)$ handle all nonseparable inputs.
Thus the preceding base factorization, infinite enumeration,
and explicit Frobenius tests supply every hypothesis of that
proved induction. They also decide validity of all further
finite tower presentations used below.

\emph{Finite certificates for places.}
Enumerate the following finite data:
\begin{enumerate}
\item elements $u,a\in K$ and a polynomial
$f(U,T)\in\mathbf F_q[U,T]$, monic in $T$ and irreducible
over $\mathbf F_q(U)$, with $f(u,a)=0$;
\item rational expressions over $\mathbf F_q$ for the original
$t$ and every supplied tower generator in $u,a$, with
nonzero denominators and with the asserted equalities in $K$;
\item a finite extension $\kappa/\mathbf F_q$ and an element
$a_0\in\kappa$ generating it, such that
$f(0,a_0)=0$ and $f_T(0,a_0)\ne0$.
\end{enumerate}
Every condition is decidable by the preceding algorithms or
finite-field computation. The second condition implies that
$K=\mathbf F_q(u,a)$. It also forces $u$ to be transcendental:
if $u$ were algebraic over $\mathbf F_q$, then the monic
equation would make $a$ algebraic, and the asserted expression
for $t$ would be impossible. Irreducibility of $f$ therefore
identifies $K$ with $\mathbf F_q(U)[T]/(f)$.

The simple root $a_0$ has a unique lift
$\alpha(U)\in\kappa[[U]]$ solving $f(U,\alpha(U))=0$.
Its coefficients are computed successively: at each degree
the correction coefficient is multiplied by the fixed unit
$f_T(0,a_0)$. Sending $a$ to $\alpha(U)$ thus gives an
injective field map $K\to\kappa((U))$.
The induced valuation has value group $\mathbf Z$, since
$u$ maps to $U$, and residue field exactly $\kappa$, since
the residue of $a$ generates $\kappa$.

For a nonzero input, substitute its rational expression in
$u,a$. The numerator and denominator are nonzero series.
Successively compute coefficients until the first nonzero
one of each appears; injectivity guarantees termination.
This gives the valuation, leading coefficient, and every
further prescribed precision by series division.
Zero inputs are first detected by exact field equality.

These embeddings have dense image. For $b\in\kappa$, express
$b=P_b(a_0)$ with $P_b\in\mathbf F_q[T]$ and use $P_b(a)$
as a lift. To approximate a finite Laurent polynomial,
correct its coefficients in increasing order, adding
$u^jP_b(a)$ at the current degree $j$. Each correction fixes
that coefficient and leaves all lower ones fixed. A finite
number of corrections attains any specified precision.
This proves both density and the asserted approximation
algorithm, and identifies the completion with
$\kappa((U))$.

\emph{Every place occurs.}
The field $K$, regardless of inseparability in its supplied
presentation, is the function field of a normal projective
curve over $\mathbf F_q$ by
\cite[Theorem 53.2.6]{C-StacksGlobalFunctionArithmetic}.
Its closed points give all places. The local rings are
discrete valuation rings and their residue fields are finite
extensions of $\mathbf F_q$.
At such a point $x$, choose a uniformizer $u\in K$ and
an affine neighbourhood on which it is regular.
The resulting finite-presentation map
$\mathbf F_q[U]\to S$, $U\mapsto u$, is \'etale at $x$:
the map on local rings is flat since it is a torsion-free
module over a discrete valuation ring, $u$ generates the
target maximal ideal, and the residue extension is finite
separable. These are exactly the hypotheses of
\cite[Lemma 10.143.7]{C-StacksGlobalFunctionArithmetic}.
After shrinking about $x$, Proposition 10.144.4 of that
source gives a standard \'etale presentation
\[
 S_g\cong \mathbf F_q[U,T]_h/(F),
 \qquad F\text{ monic in }T,
\]
with $F_T$ invertible. Let $a$ be the image of $T$ and let
$f$ be its monic minimal polynomial over $\mathbf F_q(U)$.
Because $a$ is integral over $\mathbf F_q[U]$, this monic
factor of $F$ belongs to $\mathbf F_q[U,T]$; equivalently,
use Gauss's lemma for the integrally closed polynomial ring.
Writing $F=fh$, the nonzero residue of $F_T(u,a)$ implies
that $f_T(0,a_0)\ne0$. The standard presentation gives
$K=\mathbf F_q(u,a)$ and $\kappa(x)=\mathbf F_q(a_0)$.
Consequently all the data just constructed appear in the
enumeration. Its induced valuation is the given one, since
the local presentation embeds at the chosen simple root.
This proves completeness without an oracle for normalization
or for factoring polynomials over a completed field.
\end{proof}

\begin{theorem}[Proved: splitting, division factors, and reduced norms]
\label{C:global-function-central-simple-algorithms}
Let $K$ be as in
Lemma~\ref{C:global-function-field-arithmetic} and let $A$
be a finite-dimensional central simple $K$-algebra given by
structure constants. There are ordinary algorithms which
decide whether $A$ is split, compute its index, construct
an actual division algebra $D$ and an explicit isomorphism
\[
                         A\cong M_r(D),
\]
and compute reduced norms. They also construct matrix units
and primitive idempotents. The result includes degrees
divisible by $\operatorname{char}K$ and applies uniformly
over every further finite tower.
\end{theorem}
\begin{proof}
Write $\dim_K A=n^2$. We first decide splitting by two
complementary searches.

\emph{A projective splitting variety.}
On the Grassmannian of $n$-dimensional subspaces of the
$n^2$-dimensional vector space $A$, impose stability under
left multiplication by every element of a fixed basis of
$A$. This is a projective variety $X_A$ defined by an
effectively computed finite system of homogeneous equations.
One explicit construction uses the Pl\"ucker coordinates of
$w\in\bigwedge^n A$. Impose the usual Pl\"ucker equations.
For each basis left multiplication $L_b$ and each basis
$(n-1)$-fold contraction $\iota_\xi$, also impose
\[
                     L_b(\iota_\xi w)\wedge w=0.       \tag{C.3.GA1}
\]
These are homogeneous quadratic equations, computable from
the structure constants. For a nonzero decomposable $w$,
its contractions span its associated subspace $W$, and
$z\wedge w=0$ is equivalent to $z\in W$. Thus (C.3.GA1)
is exactly left stability, in every characteristic.

For any field extension $L/K$, the points of $X_A(L)$
are exactly the $n$-dimensional left ideals of $A_L$.
If $A_L\cong M_s(D_L)$ with
$\dim_LD_L=d^2$, then $n=sd$ and every nonzero left ideal
has dimension a positive multiple of $sd^2=nd$.
The simple left module has precisely that dimension.
It follows that
\[
           X_A(L)\ne\varnothing
               \quad\Longleftrightarrow\quad A_L\text{ is split}.
                                                               \tag{C.3.GA2}
\]
These module facts follow directly from matrix-algebra
Morita equivalence; see
\cite[Lemma 11.4.6]{C-StacksGlobalFunctionArithmetic}.

\emph{Positive and negative searches.}
The positive search enumerates finite matrices over $K$
and stops when it finds a unital algebra isomorphism
$A\to M_n(K)$. All multiplication and invertibility
conditions are finite equality and linear-algebra tests.
This search terminates exactly when $A$ is split.

The negative search enumerates all place certificates of
Lemma~\ref{C:global-function-field-arithmetic} and all
positive integers $N$. For each place, multiply each
homogeneous defining equation of $X_A$ by a suitable power
of its named uniformizer so that all coefficients are
integral. Their reductions modulo $U^N$ are computable.
Exhaust all coordinate vectors over the finite ring
$\kappa_v[U]/(U^N)$ with at least one unit coordinate.
If none satisfies the equations, stop and declare $A$
nonsplit. Such a verdict is sound: a point over $K$ would
give a point over $K_v$, and a homogeneous vector representing
the latter can be scaled to integral coordinates with a
unit coordinate, giving a solution at every finite level.

Conversely, if $X_A(K_v)=\varnothing$, some finite level
has no such vector. Indeed, otherwise the finite nonempty
sets of solution vectors at levels $N$ form a finitely
branching tree under reduction, with a vertex at every
depth. An infinite path yields a compatible vector in
$\kappa_v[[U]]$ with a unit coordinate, solving all the
equations by separatedness. It would be a point of
$X_A(K_v)$, a contradiction.

If $A$ is nonsplit, the injectivity part of the
Albert--Brauer--Hasse--Noether theorem gives a place $v$
where $A_{K_v}$ is nonsplit. Hence (C.3.GA2) and the last
paragraph guarantee a terminating negative search.
We use the theorem for all global fields, including
positive characteristic, as stated and proved in
\cite[Theorem 1.5.36(i), pp.\ 25--26]{C-PoonenRationalPoints2017};
see also \cite[Section 2]{C-HurlimannABHN1992}.
Dovetailing the positive and negative searches therefore
decides splitting. No uniform bound on bad places or
congruence precision is needed.

\emph{Index and an actual division algebra.}
The period of a central simple algebra of degree $n$
divides $n$. Apply the splitting decision to tensor powers
$A^{\otimes j}$, $1\leq j\leq n$, and take the least
split exponent $d$. This is the period; it equals the
index over a global field by
\cite[Proposition 1.5.17 and Theorem 1.5.36(iv)]
{C-PoonenRationalPoints2017}. There is no restriction that
$d$ be prime to the characteristic.

Enumerate $e\in A$ until
\[
                    e^2=e,
              \qquad \dim_K Ae=nd.                    \tag{C.3.GA3}
\]
The image of right multiplication by $e$ gives $Ae$, so
both conditions are decidable. The matrix-algebra structure
theorem guarantees such an $e$. The left ideal $Ae$ is
simple because $nd$ is the minimum positive left-ideal
dimension. Consequently $D=eAe$, with identity $e$, is a
division algebra, of dimension $d^2$ over $K$.
Its multiplication table is obtained by linear algebra.
Set $r=n/d$. Enumerate a family $u_{ij}\in A$ satisfying
\[
 u_{ij}u_{k\ell}=\delta_{jk}u_{i\ell},\qquad
 \sum_i u_{ii}=1,\qquad u_{11}=e.
\]
Such a family exists for every primitive idempotent $e$
in $M_r(D)$, so the finite-equation search terminates.
It gives the explicit mutually inverse maps
\[
 A\longrightarrow M_r(D),\quad
 a\longmapsto (u_{1i}a u_{j1})_{ij},
 \qquad
 (d_{ij})\longmapsto\sum_{i,j}u_{i1}d_{ij}u_{1j}.
\]
This also produces all the diagonal primitive idempotents.

\emph{Reduced norms.}
Enumerate finite separable towers $L/K$ and finite matrices
certifying an isomorphism $A_L\cong M_n(L)$.
Every central simple algebra has a finite separable
splitting field
\cite[Proposition 11.8.5]{C-StacksGlobalFunctionArithmetic},
so this positive search terminates. For an input $a\in A$,
compute the determinant of its image. This belongs to $K$
and is independent of the chosen splitting isomorphism:
over a separable closure the isomorphisms differ by inner
automorphisms, and determinant is invariant under them.
This is precisely the reduced norm. Its value in $K$ is
read from the known tower coordinates in $L$.
The same argument computes the reduced characteristic
polynomial and applies to matrices over $A$.
The determinant descent argument is also recorded in
\cite[proof of Theorem 59.67.8]{C-StacksGlobalFunctionArithmetic}.
All constructions use the field algorithms already proved,
so they remain ordinary after any finite tower.
\end{proof}

\begin{theorem}[Proved: full finite-algebra decomposition over global function fields]
\label{C:global-function-finite-algebra-algorithms}
For every finite-dimensional associative unital algebra $E$
over a field $K$ of
Lemma~\ref{C:global-function-field-arithmetic}, its Jacobson
radical and a complete semisimple decomposition are computable:
\[
 E/\operatorname{rad}E\ \xrightarrow{\ \sim\ }
                  \prod_j M_{r_j}(D_j).
\]
The output includes the quotient map and its kernel basis,
explicit finite-tower presentations of the centers $L_j/K$,
actual central division algebras $D_j/L_j$ with multiplication
tables, mutually inverse decomposition maps, and a complete
orthogonal family of primitive idempotents in $E$ lifting
the diagonal primitive idempotents of the quotient.
The centers and supplied ground-field towers need not be
separable. Reduced norms in every simple quotient factor
are computable.
\end{theorem}
\begin{proof}
Put $h=\dim_KE$; the zero algebra is handled by the empty
product. We first find the radical and central simple
quotient factors by a positive search with finite certificates.
Enumerate finite field towers $L_j/K$, positive-dimensional
unital associative $L_j$-algebras $C_j$, and $K$-linear maps
\[
                         \pi:E\longrightarrow\prod_j C_j.
\]
The field presentations are checked by factorization,
and algebra tables are checked by finite associativity
and unit equations. For each $C_j$ require that
\[
 C_j\otimes_{L_j}C_j^{\mathrm{op}}\longrightarrow
       \operatorname{End}_{L_j}(C_j),
       \qquad a\otimes b\longmapsto(x\mapsto axb)        \tag{C.3.GA4}
\]
be an invertible linear map.

This finite test characterizes central simplicity. If it
holds, every two-sided ideal is stable under all
$L_j$-linear endomorphisms and hence is either zero or
$C_j$. A central element acts by an endomorphism commuting
with all of those endomorphisms, hence by a scalar in
$L_j$; evaluating at $1$ proves that the center is $L_j$.
Conversely, for a central simple algebra the map is an
isomorphism by
\cite[Lemma 11.4.10]{C-StacksGlobalFunctionArithmetic}.

Require further that $\pi$ preserve the unit and products,
be surjective, and have kernel $J$ with $J^{h+1}=0$.
All these are finite linear-algebra and multiplication
tests; powers of $J$ are computed as spans of products.
An accepted kernel is exactly $\operatorname{rad}E$:
a nilpotent ideal is contained in the radical, while the
quotient is a product of simple Artinian algebras and
has radical zero. The search terminates because the
ordinary Artinian structure theorem supplies the true
radical quotient, its simple factors, and their finite
field centers. Each such center, separable or not, has
a finite tower presentation, and all the required
structure constants and quotient maps occur in the
enumeration. In particular, this search does not assume
that $K$ is perfect or use a separable Wedderburn complement
inside $E$.

Apply Theorem~\ref{C:global-function-central-simple-algorithms}
to each $C_j$ over its explicit global function field $L_j$.
This produces $C_j\cong M_{r_j}(D_j)$, all actual division
tables, and their centers, reduced norms, and matrix units.
Composing with $\pi$ gives the asserted quotient map and
its explicit inverse on the quotient vector space.

Finally, complete orthogonal families of idempotents lift
modulo a nilpotent ideal. Here is the finite-algebra reason
for the existence needed by an ordinary search. If $a$ is
an approximate idempotent, its error $a^2-a$ lies in $J$;
$2a-1$ is invertible because
$(2a-1)^2=1+4(a^2-a)$. Correcting by
\[
                  a\longmapsto
               a-(2a-1)^{-1}(a^2-a)
\]
squares the error ideal, and finitely many corrections
make the error zero. To lift an orthogonal family, lift
successively in the corners complementary to the sum
already lifted, taking the final complementary idempotent.
Thus enumerate actual finite families in $E$ until their
idempotence, pairwise orthogonality, sum $1$, and prescribed
quotient images are verified. This terminates by the
lifting argument. Each lifted idempotent is primitive,
since its corner modulo the nilpotent radical is a division
algebra and a nilpotent ideal contains no nonzero idempotent.
Every requested output is therefore obtained by a
terminating ordinary algorithm.
\end{proof}
